\documentclass{article}
\usepackage{graphicx} % Required for inserting images
\usepackage[english]{babel}
\usepackage{amsmath}
\usepackage{amsthm}
\usepackage{amssymb}
\usepackage[margin=.75in]{geometry}
\usepackage{hyperref}
\usepackage{amsfonts}
\usepackage{mathrsfs}
\usepackage[version=4]{mhchem}
\usepackage{bigints}
\usepackage{caption}
\usepackage{xpatch}
\usepackage{xr}
\usepackage{enumitem}

\title{Math 381 Fourier Analysis and Boundary Value Problems for ISP Notes}
\author{Kyle Wang}
\date{September 2026}

\newcommand{\R}{\mathbb{R}}
\newcommand{\pd}[2]{\frac{\partial #1}{\partial #2}}
\newcommand{\od}[2]{\frac{d #1}{d #2}}
\newcommand{\pdt}[2]{\frac{\partial ^2 #1}{\partial #2 ^2}}
\newcommand{\odt}[2]{\frac{d^2 #1}{d #2 ^2}}

\newtheorem{theorem}{Theorem}[section]
\newtheorem{lemma}[theorem]{Lemma}
\newtheorem{corollary}[theorem]{Corollary}
\theoremstyle{definition}
\newtheorem{definition}[theorem]{Definition}
\newtheorem{example}[theorem]{Example}

\setlength{\parindent}{0pt}

\begin{document}

\maketitle
\tableofcontents
\section{Partial Differential Equations}
\begin{example}
    Find the function $u(x,y)$ such that $u: \R^2 \mapsto \R $ that satisfies the differential equation $u_{xy}=0$
\end{example}
We can just integrate twice. Once with respect to $x$ and one with respect to $y$. 
\begin{equation*}
    \begin{split}
        u_{xy} &= 0 \\
        \Rightarrow u_x&=c(x) \\
        \Rightarrow u(x,y) &=C(x) + D(y)
    \end{split}
\end{equation*}

\begin{example}
    Find a solution $u(x,y)$ to solve the transport equation 
    \begin{equation*}
        u_x+yu_y=0
    \end{equation*}
\end{example}
Conceptually we can think of this as $\vec{v}\cdot \nabla u = D_{\vec{v}}u=0$ where $\vec{v} = \langle 1,y\rangle$. We effectively want to travel along a path such that the directional derivative in $\vec{v}$ is 0. These paths are called integral curves as they are tangent to the vector field at every point. We now arrive at the characteristic equation 
\begin{equation*}
    \od{y}{x}=\frac{y}{1}
\end{equation*}
since the slope at every point is just rise over run and therefore $\frac{y}{1}$. This leads to the solution 
\begin{equation*}
    y(x) = Ae^{x}
\end{equation*}
where $A$ is found using an initial condition. We note that since the value of the function doesn't depend on where you on each curve but rather \emph{which} curve you're on, we can write $u(x,y) =f(A)$ where $A=ye^{-x}$. Therefore the general solution to the transport equation is 
\begin{equation*}
    u(x,y) = f(ye^{-x})
\end{equation*}

\end{document}
